\documentclass[12pt,a4paper]{article}

\usepackage[utf8]{inputenc}
\usepackage[T1]{fontenc}
\usepackage{lmodern}
\usepackage[british]{babel}
\usepackage{amsmath}
\usepackage[margin=2cm]{geometry}
\usepackage{setspace}
\usepackage[hidelinks]{hyperref}

\onehalfspacing
\setlength{\parindent}{0pt}
\setlength{\parskip}{0.4em}

\hypersetup{
  pdftitle={Research Proposal: Carbon Leakage through the Back Door?},
  pdfauthor={Baha Khammari}
}

\begin{document}

\begin{center}
{\LARGE\bfseries Proposal Master's Thesis}\\[0.2cm]
{\large\bfseries University of Cologne -- Chair in Economics, Energy and Sustainability}
\end{center}

\vspace{0.5cm}

\begin{tabular}{@{}ll@{}}
Name: & Baha \\
Surname: & Khammari \\
Student ID Number: & 7364599 \\
Course of studies: & MSc Economics \\
Semester: & Summer Term 2026 \\
\end{tabular}

\vspace{0.6cm}

\subsection*{Title}

Carbon Leakage through the Back Door? The Effective EU~ETS Carbon Cost and EU Imports of Emission-Intensive Goods from North Africa

\subsection*{Research Question (one sentence)}

Did increases in the effective EU~ETS carbon cost, net of free allocation, raise EU imports of emission-intensive goods from the North African partners Morocco, Algeria, Egypt, and Tunisia, relative both to less emission-intensive goods and to comparison partners without carbon pricing, over 2005 to 2021?

\section{Motivation and background of the research question (max.\ 250 words)}

The EU~ETS prices carbon for European producers while most trading partners price it weakly or not at all. Theory predicts that this asymmetry shifts sourcing of emission-intensive goods towards lower-regulation partners, so that foreign emissions rise and partly offset domestic abatement. The empirical record sits uneasily with this prediction: aggregate studies of European trade find no evidence of leakage (Naegele and Zaklan, 2019; Branger et al., 2016), and firm-level studies find emissions reductions without competitiveness losses (Colmer et al., 2025). Aggregate estimates, however, can conceal effects concentrated in particular bilateral corridors, and most samples end before the allowance price became economically meaningful after 2018.

The thesis would fill this gap by estimating leakage where it is most likely to appear: the North African corridor, whose four partners combine proximity, existing cement, fertiliser, and steel capacity, preferential access under Euro-Mediterranean agreements (European Commission, 2026), and no carbon pricing of their own within the sample period (World Bank, 2026). The question has direct policy stakes. The Carbon Border Adjustment Mechanism entered its definitive regime in January 2026, and its calibration presumes that carbon-cost differentials divert trade. A corridor-level estimate of the pre-CBAM baseline would inform this calibration and provide the empirical counterpart to simulation-based evaluations (Bellora and Fontagn\'{e}, 2023). Because the price acceleration begins in 2019 and the pandemic years are set aside, the clean strong-treatment window is short, so the thesis reports effect sizes with their precision rather than resting on significance.

\section{Methodology and Data (max.\ 250 words)}

The thesis would estimate a triple difference on EU import flows, 2000 to 2021, interacting exposure $\tau_p$, carbon cost $C_{st}$, and corridor indicator $L_c$, with product-by-partner, partner-by-year, and product-by-year fixed effects. The comparison group is the Western Balkans, which priced no carbon. Because $C_{st}$ varies across sectors within a partner-year cell, $C_{st} \times L_c$ is not absorbed and enters as a control.

The treatment $C_{st} = P_t \times (1 - a_{st})$ nets the free-allocation share, measured against fixed 2005--2007 emissions and capped at one, out of the EUA price. This presumes allocation lowers marginal cost, which holds because activity thresholds and closure rules tie allowances to production (European Commission, 2011). Exposure $\tau_p$ is sectoral CO$_2$ intensity at NACE Rev.~2, fixed at 2008--2010, in tonnes per euro of output, so $\tau_p C_{st}$ is a carbon-cost share and the coefficient a percentage change per point.

Data are public: Eurostat Comext (DS-045409), the EU Transaction Log, requiring a NACE crosswalk, and Eurostat air emissions accounts, available from 2008. Both value and quantity are estimated, since carbon costs raise unit values. Estimation uses Poisson pseudo-maximum likelihood, correcting heteroskedasticity bias in log-linear models (Santos Silva and Tenreyro, 2006). With a continuous dose, two-way fixed effects can weight comparisons non-convexly (de Chaisemartin and D'Haultf\oe uille, 2020; Callaway et al., 2024), so dose-specific responses are reported.

Parallel trends is the central assumption, tested over 2000 to 2002. Four threats follow: four treated clusters requiring randomisation inference, Algerian and Egyptian gas price regulation, the 2019 EU steel safeguard, and the pandemic years.

\section{Literature Review (max.\ 250 words)}

Trade-flow studies find little evidence of leakage: Naegele and Zaklan (2019) measure the EU~ETS emission cost net of free allocation in GTAP embodied-carbon flows and find none, while Branger et al.\ (2016) find no carbon-price effect on EU net imports of cement and steel. Aichele and Felbermayr (2015) provide the contrasting result that binding Kyoto commitment raised embodied-carbon imports from non-committed partners by 8 per cent. Firm-level evidence agrees: Colmer et al.\ (2025) find emissions reductions of 14 per cent in Phase~I and 16.3 per cent in Phase~II for French manufacturing firms with no loss of value added or employment and no evidence of outsourcing, Dechezlepr\^{e}tre et al.\ (2023) find reductions of about 10 per cent alongside revenue gains, and Koch and Basse Mama (2019) find investment leakage confined to footloose, less capital-intensive German firms. Surveys attribute these nulls to generous free allocation and low allowance prices in the periods sampled (Joltreau and Sommerfeld, 2019; Verde, 2020). Ex-ante simulations nevertheless predict substantial leakage at higher carbon prices (Bellora and Fontagn\'{e}, 2023).

Two features distinguish the proposed thesis from the closest existing study. Naegele and Zaklan (2019) enter importer and exporter stringency as separate terms, never interacted with partner regulatory status, so an effect confined to one corridor is not tested. Their GTAP 9.2 benchmark years are 2004, 2007 and 2011, so one pre-policy observation carries the counterfactual. The proposed design differentiates partners, runs annually through 2021, and has three pre-policy years.

\section{Table of Contents (including page numbers)}

\begin{tabbing}
\hspace{1.2cm} \= \hspace{12.5cm} \= \kill
1 \> Introduction \> 1 \\
2 \> Policy Background: the EU~ETS, Free Allocation, and the CBAM \> 3 \\
3 \> Literature Review \> 7 \\
4 \> Empirical Strategy \> 12 \\
\> \hspace{0.5cm} 4.1 Specification and Identification \> 12 \\
\> \hspace{0.5cm} 4.2 Threats and Identifying Assumptions \> 15 \\
5 \> Data and Variable Construction \> 18 \\
6 \> Results \> 23 \\
7 \> Robustness and Heterogeneity \> 30 \\
8 \> Discussion: Implications for the CBAM \> 37 \\
9 \> Conclusion \> 40 \\
\> References \> 42 \\
\> Appendix: Additional Tables and Event-Study Figures \> 46 \\
\end{tabbing}

Planned length: approximately 50 pages including appendix, within the 60-page limit for a master's thesis.

\section{References (max.\ 500 words)}

\begin{spacing}{1.0}

\textbf{Aichele, R. and Felbermayr, G. (2015):} Kyoto and Carbon Leakage: An Empirical Analysis of the Carbon Content of Bilateral Trade. In: Review of Economics and Statistics, 97(1), pp.\ 104--115.

\textbf{Bellora, C. and Fontagn\'{e}, L. (2023):} EU in Search of a Carbon Border Adjustment Mechanism. In: Energy Economics, 123, 106673.

\textbf{Branger, F., Quirion, P. and Chevallier, J. (2016):} Carbon Leakage and Competitiveness of Cement and Steel Industries under the EU ETS: Much Ado about Nothing. In: The Energy Journal, 37(3), pp.\ 109--136.

\textbf{Callaway, B., Goodman-Bacon, A. and Sant'Anna, P.\ H.\ C. (2024):} Event Studies with a Continuous Treatment. In: AEA Papers and Proceedings, 114, pp.\ 601--605.

\textbf{de Chaisemartin, C. and D'Haultf\oe uille, X. (2020):} Two-Way Fixed Effects Estimators with Heterogeneous Treatment Effects. In: American Economic Review, 110(9), pp.\ 2964--2996.

\textbf{Colmer, J., Martin, R., Mu\^{u}ls, M. and Wagner, U.\ J. (2025):} Does Pricing Carbon Mitigate Climate Change? Firm-Level Evidence from the European Union Emissions Trading System. In: Review of Economic Studies, 92(3), pp.\ 1625--1660.

\textbf{Dechezlepr\^{e}tre, A., Nachtigall, D. and Venmans, F. (2023):} The Joint Impact of the European Union Emissions Trading System on Carbon Emissions and Economic Performance. In: Journal of Environmental Economics and Management, 118, 102758.

\textbf{European Commission (2011):} Decision 2011/278/EU Determining Transitional Union-Wide Rules for Harmonised Free Allocation of Emission Allowances. Official Journal of the European Union, L 130/1, 17.05.2011.

\textbf{European Commission (2026):} EU Trade Relations with the Southern Neighbourhood. Directorate-General for Trade. \url{https://policy.trade.ec.europa.eu/eu-trade-relationships-country-and-region/countries-and-regions/southern-neighbourhood_en}. Accessed: 05.08.2026.

\textbf{Joltreau, E. and Sommerfeld, K. (2019):} Why Does Emissions Trading under the EU Emissions Trading System Not Affect Firms' Competitiveness? Empirical Findings from the Literature. In: Climate Policy, 19(4), pp.\ 453--471.

\textbf{Koch, N. and Basse Mama, H. (2019):} Does the EU Emissions Trading System Induce Investment Leakage? Evidence from German Multinational Firms. In: Energy Economics, 81, pp.\ 479--492.

\textbf{Naegele, H. and Zaklan, A. (2019):} Does the EU ETS Cause Carbon Leakage in European Manufacturing? In: Journal of Environmental Economics and Management, 93, pp.\ 125--147.

\textbf{Santos Silva, J.\ M.\ C. and Tenreyro, S. (2006):} The Log of Gravity. In: Review of Economics and Statistics, 88(4), pp.\ 641--658.

\textbf{Verde, S.\ F. (2020):} The Impact of the EU Emissions Trading System on Competitiveness and Carbon Leakage: The Econometric Evidence. In: Journal of Economic Surveys, 34(2), pp.\ 320--343.

\textbf{World Bank (2026):} Carbon Pricing Dashboard. \url{https://carbonpricingdashboard.worldbank.org}. Accessed: 05.08.2026.

\end{spacing}

\end{document}
